% Status: DRAFT (revised per Reviewer-2)
\documentclass[11pt]{article}
\usepackage[margin=1in]{geometry}
\usepackage{booktabs}
\usepackage{amsmath}
\usepackage{url}
\usepackage{microtype}

\title{Auditing Evaluation Pipelines with Cheap Models:\\ An Unsupervised Miner for Verdict Errors}
\author{Anonymous}
\date{August 2026 (draft)}

\begin{document}
\maketitle

\begin{abstract}
Automatic benchmark evaluations produce a per-record 0/1 correctness verdict from an answer key plus a scorer, and downstream work treats that verdict as ground truth.
When models disagree with the verdict, the standard reading is that the models failed.
We show that for one specific and cheap signal the opposite reading is usually correct.
We flag any record where three or more solver models, all officially graded incorrect, nevertheless agree with one another on the final answer.
On GSM8K, validated externally against GSM8K-Platinum human re-annotation~\cite{platinum} plus mechanical checks, 89\% of these flags (66 of 74) mark genuine evaluation-pipeline verdict errors, decomposing into 72\% bugs in our own scorer, 15\% items Platinum removed as broken, and roughly 3\% genuinely wrong answer keys.
Measured against dataset-level defects alone (broken items plus wrong keys), the same flags have 18\% precision (13 of 74).
A second, more expensive signal, in which two frontier models independently re-solve the problem and both contradict the official verdict, reaches 97\% precision (39 of 40) on its flagged slice, subject to an adjudicator-circularity caveat.
The miner is unsupervised and needs no human labels to run.
Our motivation is cautionary: three apparent judge-failure laws in our own experiments dissolved under adjudication; in the attractor case study, 96\% of flagged cases (83\% scorer bugs plus 13\% label errors or suspect items) were not judge failures, and in the detection-hierarchy case study, 71\% were verdict errors.
We release the audit and a ledger of 751 correctness-label upgrades across GSM8K, GSM8K-Platinum, and a Putnam-Axiom variant.
The implication for the judge-evaluation literature is direct: scoring judges against noisy verdicts can manufacture failure laws, so pipeline audits must come first.
\end{abstract}

\section{Introduction}

Every automatic benchmark evaluation rests on a chain of trust.
The dataset supplies an answer key, a scorer compares model output to that key, and the resulting 0/1 verdict is treated as ground truth.
Work on LLM judges inherits this chain.
A judge that calls an officially-wrong solution correct is recorded as a false positive, and aggregate false-positive rates become claims about judge reliability.

This paper reports what happened when we stopped trusting the chain.
In a study of judge reliability on GSM8K and MATH we found three clean, quantitative laws of judge failure.
All three dissolved under adjudication.
In the attractor case study, 96\% of the flagged cases (83\% scorer bugs plus 13\% label errors or suspect items) were not judge failures at all.
In the detection-hierarchy case study, 71\% of the flagged cases were verdict errors.
The failures were wrong answer keys, broken questions, or bugs in our own string-matching scorer, not failures of the judges.

The dissolution suggested an inversion.
If disagreement with the official verdict is dominated by verdict noise, then such disagreement is not a model-failure detector.
It is a verdict-error detector.
We operationalize this as a miner whose primary signal requires no judge at all: flag any record where three or more solver models, all officially graded incorrect by the pipeline, agree with one another on the extracted final answer.
Validated against GSM8K-Platinum human re-annotation~\cite{platinum} and mechanical checks, and using no frontier model anywhere in flagging or adjudication, 89\% of the flags (66/74) mark genuine pipeline-verdict errors, and 18\% (13/74) mark dataset-level defects.

Our contributions are:
(1) three documented case studies in which plausible judge-failure laws reduce to verdict artifacts;
(2) an unsupervised verdict-error miner with externally validated 89\% precision (66/74) for pipeline-verdict errors and 18\% (13/74) for dataset-level defects, plus a frontier-arbitrated variant at 97\% (39/40) that carries a circularity caveat;
(3) released artifacts, including a 751-record ledger of correctness-label corrections.

\section{The Failure Mode: Three Laws That Dissolved}
\label{sec:laws}

All three case studies come from a study of judge reliability on GSM8K and MATH-500 solutions produced by 17 open-weights solvers, scored by an automatic pipeline (answer key plus string-matching scorer), and judged by a ladder of Qwen2.5-Instruct models (0.5B--14B) plus two frontier judges (GPT-5.2 and Claude Opus).
Each judge scores a candidate solution on a 0--10 scale, where 0 means certainly incorrect and 10 means certainly correct.
We call a score of at least 7 on an officially-wrong solution a judge false positive.
Table~\ref{tab:laws} summarizes the pattern.

\begin{table}[t]
\centering
\small
\begin{tabular}{@{}p{3.1cm}p{4.6cm}p{5.6cm}@{}}
\toprule
Apparent law & Raw signal & After adjudication \\
\midrule
Attractor effect & Judges fooled by ``attractor'' wrong answers; within-item score gap $+6.41$ on GSM8K & 83\% of the 101 flagged records were scorer bugs and 13\% label errors or suspect items (96\% non-judge causes); on the 4\% genuine residue, judge false positives were 0\% \\
\addlinespace
Capability-gap degradation & Judge false positives peak on stronger solvers' errors & Frontier-adjudicated residual disagreement is $\sim$2\% (circularity caveat, Section~\ref{sec:capgap}) and flat across solver strength; universal fools are $\sim$0.2\% of wrong solutions \\
\addlinespace
Detection hierarchy & $P(\text{frontier fooled})$ rises $1\%\!\to\!6.5\%\!\to\!60.5\%\!\to\!94\%$ as $k=0..3$ weak judges are fooled & 71\% of $k\ge2$ cases were verdict errors; conditioned on confirmed-wrong, frontier FP at $k\ge2$ is 6.2\% (GPT-5.2) and 8.3\% (Opus) vs.\ baselines of 2.0\% and 1.5\% \\
\bottomrule
\end{tabular}
\caption{Three judge-failure laws and what adjudication left of them.}
\label{tab:laws}
\end{table}

\subsection{The attractor effect}

The first pass suggested that wrong answers shared by many solvers (``attractors'') systematically fool judges.
On GSM8K, attractor wrong answers received judge scores 6.41 points higher, on the 0--10 scale, than other wrong answers to the same item.
Adjudicating all 101 attractor records against robust answer normalization and GSM8K-Platinum dissolved the effect.
83\% were bugs in our own scorer: the model's answer equaled the gold answer, the string matcher misparsed it, and both judges had scored the solution near 9.7 out of 10.
The judges were right and the oracle was wrong.
Another 13\% were label errors or suspect items, so 96\% of the flagged cases had non-judge causes.
On the 4\% of records that were genuinely wrong, judges scored a mean of 2.5 with a 0\% false-positive rate.

\subsection{The capability-gap confound}
\label{sec:capgap}

The second pass suggested that judge reliability degrades when solvers outgrow judges, the folklore version of the oversight problem.
The raw judge-by-solver surface even showed a false-positive peak at the 14B-solver column for every judge.
After flipping the adjudicated verdict errors, the frontier false-positive rate on genuine natural errors is about 2\%, flat across solver strength, weak-judge consensus, and difficulty.
This $\sim$2\% figure is a residual of frontier-adjudicated disagreement rather than a ground-truth error rate, because correlated frontier families supply both the judgments and the adjudication.
Errors that fool every judge are about 0.2\% of wrong solutions.
On cleaned wrong sets, frontier false positives (score $\ge 7$) are 2--3\% on MATH and 0.7--0.9\% on Putnam.
The GSM8K residual of $\sim$17\% still contained unflipped verdict errors at the time of measurement.

\subsection{The detection hierarchy}

The third pass produced the sharpest raw pattern.
The number of weak judges (1.5B/3B/7B) fooled by a wrong candidate seemed to predict frontier failure: $P(\text{frontier fooled})$ rose from 1\% at $k=0$ to 94\% at $k=3$, a $\sim$94$\times$ lift.
Adjudication dissolved it.
A solve-first binary adjudication of all 55 cases with $k\ge2$ found that 71\% were verdict errors: the candidate solution was actually correct (GSM8K 32/40 even after Platinum cleaning; MATH 7/15).
An independent arbiter from a second frontier family agreed on 55 of 55 adjudications.
Conditioned on confirmed-wrong solutions, frontier false positives at $k\ge2$ are 6.2\% for GPT-5.2 and 8.3\% for Opus, against baselines of 2.0\% and 1.5\% respectively.
Weak-judge consensus carries almost no signal about frontier failure on genuine errors.
What it detects, with high reliability, is bad verdicts.

\section{The Miner}
\label{sec:miner}

The dissolved patterns are not subtlety meters.
They are verdict-error miners, and we use them as such.
Two distinct signals are involved, and we describe them separately because they were generated by different model pools, validated by different adjudicators, and achieve different precisions.

\paragraph{Signal 1: multi-solver answer consensus (primary; source of the 74 flags).}
For each benchmark record (a problem, candidate solutions, and the official 0/1 correctness verdicts from key plus scorer), collect the extracted final answers of the solver pool used in the evaluation.
Flag the item when three or more solvers whose solutions the pipeline graded incorrect nevertheless agree with one another on the final answer.
The consensus answer may equal the official key, which indicates a scorer misparse, or differ from it, which indicates a wrong key or a broken item.
The consensus pools in the released audit span three to nine solvers per flag and include models from the Mistral, Microsoft Phi, DeepSeek, AI2 OLMo, and Qwen families.
These are the evaluation's solvers, not judges: many of them (e.g.\ Mathstral, Phi-3.5-mini, deepseek-math-7b, OLMo-2, Ministral) appear in no judge ladder.
The signal is unsupervised, involves no threshold tuning on labels and no human in the loop, and adds no model calls beyond the solver runs the evaluation already performed.
This signal generated all 74 flags validated in Section~\ref{sec:validation} at 89\% precision for pipeline-verdict errors.

\paragraph{Signal 2: frontier re-solve verdicts (separate, more expensive).}
Independently of Signal 1, two frontier models (GPT-5.2 and Claude Opus) each re-solve the problem from scratch and issue a binary verdict on the candidate solution.
A record is flagged when both re-solve verdicts contradict the official verdict.
In the measured cell this signal reached 97\% precision (39 of 40), and the two frontier families agreed on 55 of 55 adjudications.
This figure is frontier-adjudicated: correlated frontier families both raise and adjudicate the flags, so the 97\% carries an adjudicator-circularity caveat and we do not treat it as externally validated.

\paragraph{A third, related variant: judge-ladder unanimity.}
A ladder of small open-weights judges (Qwen2.5-Instruct at 1.5B, 3B, and 7B) can also be run over each record, flagging when every judge on the ladder disagrees with the official verdict.
A compliance gate (parseable verdicts plus planted calibration pairs) screens out judges that cannot follow the protocol; specialized models can fail this gate.
This variant is related to Signals 1 and 2 but is not what generated the 74 externally validated flags, and the numbers in Section~\ref{sec:validation} should not be attributed to it.

\paragraph{What a flag means.}
A flag does not distinguish among failure types by itself.
Adjudication of flags separates three error classes: wrong answer keys, broken or ambiguous questions, and scorer bugs (the model's answer matches the key but the grading harness misparses it).
All three corrupt the recorded verdict, and therefore corrupt downstream judge evaluation, equally.

\section{Validation}
\label{sec:validation}

\paragraph{External validation, no frontier models.}
The circularity risk in Section~\ref{sec:laws} is that frontier models adjudicated flags raised partly by frontier models.
We therefore validated the consensus miner (Signal 1) externally.
We ran multi-solver consensus flags on original GSM8K verdicts and adjudicated every flag using only GSM8K-Platinum human re-annotation~\cite{platinum} and mechanical answer checks, with no frontier model anywhere in flagging or adjudication.
Result: 66 of 74 flags were genuine pipeline-verdict errors, a precision of 89\% (Wilson 95\% interval $[0.80, 0.94]$).
Measured against dataset-level defects alone (broken items plus wrong keys), precision is 18\% (13 of 74).
Table~\ref{tab:validation} gives the breakdown.\footnote{The counts here differ superficially from the audit file, which lists 60 scorer bugs and 4 Platinum-removed items: 7 items appear in the audit's scorer-bug appendix but were also removed by Platinum, and the table counts each flag once, assigning those 7 double-counted items to the broken-item row ($53 = 60 - 7$; $11 = 4 + 7$).}

\begin{table}[t]
\centering
\begin{tabular}{@{}lrr@{}}
\toprule
Flag outcome & Count & Share \\
\midrule
Scorer bug (answer matches key; grader misparsed) & 53 & 72\% \\
Broken item (removed by Platinum re-annotation) & 11 & 15\% \\
Wrong key (relabeled by Platinum) & 2 & 3\% \\
\midrule
Confirmed pipeline-verdict error (total) & 66 & 89\% \\
\quad of which dataset-level defect (broken $+$ wrong key) & 13 & 18\% \\
Models wrong together (miner false alarm) & 8 & 11\% \\
\midrule
Total flags & 74 & 100\% \\
\bottomrule
\end{tabular}
\caption{External adjudication of multi-solver consensus flags on GSM8K against GSM8K-Platinum and mechanical checks. No frontier model is used in flagging or adjudication.}
\label{tab:validation}
\end{table}

\paragraph{Frontier-arbitrated operating point.}
With the separate frontier re-solve signal (Signal 2), precision on the flagged slice was 97\% (39/40), and the two frontier families agreed on 55/55 adjudications.
We report 89\% as the headline because it is the externally validated number; the 97\% figure inherits adjudicator circularity (Section~\ref{sec:miner}).

\paragraph{Seed audit examples.}
The qualitative audit behind these counts is released alongside the paper.
Two examples of genuine key errors on GSM8K, both confirmed by Platinum:
item 403 (air-conditioner energy) has official key 81 while the model consensus and Platinum give 135;
item 454 (apples eaten by two people) has official key 150 while consensus and Platinum give 240.
The audit also lists 4 items Platinum removed as broken or ambiguous, a further 7 removed items recorded in the scorer-bug appendix, 8 contrast cases where the flag fired but the key was right, and 60 scorer bugs in our own grading pipeline.

\paragraph{Effect on measured judge quality.}
Verdict errors are not cosmetic.
In our own pipeline the scorer-bug rate over the GSM8K records in the GPT-judged runs was 4.1\% overall, and cleaning raised the measured judge-quality signal on that slice from 0.656 to 0.753.
Noisy verdicts had been understating judge quality throughout.

\section{Released Artifacts}

\paragraph{Correction ledger.}
We release \texttt{cc\_corrections.jsonl}, a ledger of 751 correctness-label upgrades (0 to 1) produced by re-scoring legacy score files with robust answer normalization backed by a math-equivalence checker.
The upgrades concentrate exactly in the diagnosed bug class: 256 on GSM8K, 249 on GSM8K-Platinum, 246 on a Putnam-Axiom variant, and zero on MATH files that use boxed-answer formatting.
The retrofit is upgrade-only: records originally scored correct are never re-litigated, because the validated failure mode of the scorer is false negatives.
Spot-checked upgrades are formatting mismatches such as \texttt{\$57.00} versus \texttt{57}, ``Answer: 38'', and comma-grouped numbers.

\paragraph{Audit file.}
We release the full seed audit of flagged GSM8K and MATH-500 items, split into genuine key errors, Platinum-removed items, frontier-adjudicated verdict errors (39 of 55 matrix-study cases), contrast cases where the key was right, and scorer bugs.

\paragraph{Miner recipe.}
The consensus rule, the frontier re-solve protocol, and the judge-ladder variant (ladder, prompts, compliance gate, unanimity rule) are specified so that the miner can be re-run on any benchmark with per-record correctness verdicts.

\section{Limitations}

\paragraph{The target is the recorded verdict, not only the key.}
The 89\% precision is for the recorded correctness verdict, which is the product of key and scorer.
72\% of confirmed errors were scorer bugs in our own harness.
Precision for dataset-level defects alone is 18\% (13/74: 15\% broken items plus 3\% relabels).
Teams with more robust scorers should expect lower flag yields and a higher share of dataset defects among flags.

\paragraph{Domain scope.}
All validation is on oracle-verifiable math (GSM8K, with supporting evidence from MATH-500).
Precision on other domains, including code and open-ended tasks, is untested here.

\paragraph{Single benchmark for the headline number.}
The external 89\% figure comes from one benchmark and 74 flags.
The Wilson 95\% interval for 66/74 is $[0.80, 0.94]$.

\paragraph{Recall.}
We validate precision on the flagged slice.
Because GSM8K-Platinum is an exhaustive re-annotation of the GSM8K test set, dataset-defect recall is computable as the fraction of Platinum-removed or relabeled items in our evaluated slice that the miner flagged:
\[
\text{recall} \;=\; \frac{\lvert \text{flagged} \cap \text{Platinum-removed/relabeled} \rvert}{\lvert \text{Platinum-removed/relabeled items in the evaluated slice} \rvert}.
\]
The numerator on our slice is 13 (11 removed plus 2 relabeled).
Computing it: 46 of the 110 Platinum-removed items fall inside our evaluated 500-item slice (question-text diff between GSM8K and GSM8K-Platinum test splits), and the miner flagged 13 of the 46, a dataset-defect recall of $28\%$.
Low recall is expected by construction: the consensus signal fires only when several solvers agree on one alternative answer, which most defective items do not elicit.
The miner is a precision instrument, not a sweep.
Recall over scorer bugs is separately unmeasured, since no exhaustive ground-truth audit of the full slice exists on our side.

\paragraph{Adjudicator dependence at the 97\% point.}
The frontier-arbitrated figure inherits adjudicator circularity, since correlated frontier families both flag and adjudicate.
That is exactly why we report the Platinum-only 89\% figure as primary.

\section{Related Work}

\paragraph{Benchmark label noise.}
Label errors in widely-used test sets are a documented, general phenomenon~\cite{pervasive}, and confident-learning methods estimate and clean such noise from model confidence statistics~\cite{confident}.
GSM8K-Platinum re-annotates GSM8K~\cite{gsm8k} with human review and removes or relabels defective items~\cite{platinum}.
Our delta is not the existence of label noise but the instrument and the audit: answer consensus among the open-weights solvers an evaluation already runs suffices to mine pipeline-verdict errors unsupervised at 89\% precision, where prior cleaning pipelines used frontier models and humans.
Our audit also shows that most mined errors (72\%) live in the evaluation pipeline's own scorer rather than in the dataset, a layer that dataset-cleaning work does not target.

\paragraph{Judge evaluation.}
The judge-evaluation literature routinely scores judges against benchmark oracles, from holistic LLM-as-a-judge protocols such as MT-Bench~\cite{mtbench} to step-level error identification in ProcessBench~\cite{processbench}, on datasets including GSM8K and MATH~\cite{gsm8k,math}.
Our case studies show that this practice can manufacture quantitative judge-failure laws from verdict noise.
We view the miner as a mandatory preprocessing step for that literature: run it, adjudicate the flags, then measure the judge.

\begin{thebibliography}{9}

\bibitem{gsm8k}
K.~Cobbe et al.
Training Verifiers to Solve Math Word Problems.
arXiv:2110.14168, 2021.

\bibitem{math}
D.~Hendrycks et al.
Measuring Mathematical Problem Solving With the MATH Dataset.
arXiv:2103.03874, 2021.

\bibitem{platinum}
J.~Vendrow, E.~Vendrow, S.~Beery, and A.~Madry.
Do Large Language Model Benchmarks Test Reliability?
arXiv:2502.03461, 2025.
(Introduces the platinum benchmarks, including GSM8K-Platinum.)

\bibitem{pervasive}
C.~G. Northcutt, A.~Athalye, and J.~Mueller.
Pervasive Label Errors in Test Sets Destabilize Machine Learning Benchmarks.
arXiv:2103.14749, 2021.

\bibitem{confident}
C.~G. Northcutt, L.~Jiang, and I.~L. Chuang.
Confident Learning: Estimating Uncertainty in Dataset Labels.
arXiv:1911.00068, 2019.

\bibitem{mtbench}
L.~Zheng et al.
Judging LLM-as-a-Judge with MT-Bench and Chatbot Arena.
arXiv:2306.05685, 2023.

\bibitem{processbench}
C.~Zheng et al.
ProcessBench: Identifying Process Errors in Mathematical Reasoning.
arXiv:2412.06559, 2024.

\end{thebibliography}

\end{document}
